\documentclass[10pt,a4paper]{amsart}
\usepackage{amsmath,amssymb,amsthm,graphicx}
% Retained preamble for the Stacks Project chapter divisors.tex (Divisors).
% The chapter file pinned-original-source.tex carries only \input{preamble} on its first line, so no
% preamble text lives inside the chapter itself.  The Stacks Project uses one single shared preamble for
% every chapter; the snapshot stages it verbatim in sources/stacks-chow/preamble.tex.  The block below is
% copied mechanically, line for line, from that file's lines 14-18 (the Xy-pic setup) and lines 154-217
% (the native theorem declarations together with the mathematical macros).  Nothing was added, removed or
% redefined.  The wrapper class and the amsmath/amssymb/amsthm/graphicx base packages are supplied by the
% assembler, so the class/\externaldocument part of the shared preamble is not repeated here.
% For commutative diagrams we use Xy-pic
\usepackage[all]{xy}

% We use 2cell for 2-commutative diagrams.
\xyoption{2cell}
% Theorem environments.
%
\theoremstyle{plain}
\newtheorem{theorem}[subsection]{Theorem}
\newtheorem{proposition}[subsection]{Proposition}
\newtheorem{lemma}[subsection]{Lemma}

\theoremstyle{definition}
\newtheorem{definition}[subsection]{Definition}
\newtheorem{example}[subsection]{Example}
\newtheorem{exercise}[subsection]{Exercise}
\newtheorem{situation}[subsection]{Situation}

\theoremstyle{remark}
\newtheorem{remark}[subsection]{Remark}
\newtheorem{remarks}[subsection]{Remarks}

\numberwithin{equation}{subsection}

% Macros
%
\def\lim{\mathop{\mathrm{lim}}\nolimits}
\def\colim{\mathop{\mathrm{colim}}\nolimits}
\def\Spec{\mathop{\mathrm{Spec}}}
\def\Hom{\mathop{\mathrm{Hom}}\nolimits}
\def\Ext{\mathop{\mathrm{Ext}}\nolimits}
\def\SheafHom{\mathop{\mathcal{H}\!\mathit{om}}\nolimits}
\def\SheafExt{\mathop{\mathcal{E}\!\mathit{xt}}\nolimits}
\def\Sch{\mathit{Sch}}
\def\Mor{\mathop{\mathrm{Mor}}\nolimits}
\def\Ob{\mathop{\mathrm{Ob}}\nolimits}
\def\Sh{\mathop{\mathit{Sh}}\nolimits}
\def\NL{\mathop{N\!L}\nolimits}
\def\CH{\mathop{\mathrm{CH}}\nolimits}
\def\proetale{{pro\text{-}\acute{e}tale}}
\def\etale{{\acute{e}tale}}
\def\QCoh{\mathit{QCoh}}
\def\Ker{\mathop{\mathrm{Ker}}}
\def\Im{\mathop{\mathrm{Im}}}
\def\Coker{\mathop{\mathrm{Coker}}}
\def\Coim{\mathop{\mathrm{Coim}}}

% Boxtimes
%
\DeclareMathSymbol{\boxtimes}{\mathbin}{AMSa}{"02}

%
% Macros for moduli stacks/spaces
%
\def\QCohstack{\mathcal{QC}\!\mathit{oh}}
\def\Cohstack{\mathcal{C}\!\mathit{oh}}
\def\Spacesstack{\mathcal{S}\!\mathit{paces}}
\def\Quotfunctor{\mathrm{Quot}}
\def\Hilbfunctor{\mathrm{Hilb}}
\def\Curvesstack{\mathcal{C}\!\mathit{urves}}
\def\Polarizedstack{\mathcal{P}\!\mathit{olarized}}
\def\Complexesstack{\mathcal{C}\!\mathit{omplexes}}
% \Pic is the operator that assigns to X its picard group, usage \Pic(X)
% \Picardstack_{X/B} denotes the Picard stack of X over B
% \Picardfunctor_{X/B} denotes the Picard functor of X over B
\def\Pic{\mathop{\mathrm{Pic}}\nolimits}
\def\Picardstack{\mathcal{P}\!\mathit{ic}}
\def\Picardfunctor{\mathrm{Pic}}
\def\Deformationcategory{\mathcal{D}\!\mathit{ef}}

% --- appended 2026-10-12 (Chief): the project own definitions of the reference/slogan/history
% display environments, copied mechanically from the shared Stacks preamble lines 8-12
% (\usepackage{verbatim} plus the three \newenvironment lines, which the chapter preamble keeps
% verbatim from the project preamble). Nothing added, removed or redefined: these environments
% carry no mathematical content and the project itself typesets them as comments. Without these
% three lines a counted statement whose own text uses them fails with Environment slogan undefined.
\usepackage{verbatim}
\newenvironment{reference}{\comment}{\endcomment}
\newenvironment{slogan}{\comment}{\endcomment}
\newenvironment{history}{\comment}{\endcomment}
\usepackage{cleveref}
\crefname{originallemma}{original lemma}{original lemmas}
\Crefname{originallemma}{Original lemma}{Original lemmas}
\makeatletter
\newlabel{lemma-base-change-fitting-ideal}{{0C3D}{source}{}{sourcepointer-1}{}}
\newlabel{lemma-base-change-fitting-ideal@cref}{{[originallemma][0][]0C3D}{[1][0][]source}{}{}{}}
\makeatother
\crefname{originallemma}{original lemma}{original lemmas}
\Crefname{originallemma}{Original lemma}{Original lemmas}
\makeatletter
\newlabel{lemma-fitting-ideal-finite-locally-free}{{0C3G}{source}{}{sourcepointer-2}{}}
\newlabel{lemma-fitting-ideal-finite-locally-free@cref}{{[originallemma][0][]0C3G}{[1][0][]source}{}{}{}}
\makeatother
\crefname{originallemma}{original lemma}{original lemmas}
\Crefname{originallemma}{Original lemma}{Original lemmas}
\makeatletter
\newlabel{lemma-fitting-ideal-generate-locally}{{0C3F}{source}{}{sourcepointer-3}{}}
\newlabel{lemma-fitting-ideal-generate-locally@cref}{{[originallemma][0][]0C3F}{[1][0][]source}{}{}{}}
\makeatother
\crefname{originallemma}{original lemma}{original lemmas}
\Crefname{originallemma}{Original lemma}{Original lemmas}
\makeatletter
\newlabel{lemma-fitting-ideal-of-finitely-presented}{{0C3E}{source}{}{sourcepointer-4}{}}
\newlabel{lemma-fitting-ideal-of-finitely-presented@cref}{{[originallemma][0][]0C3E}{[1][0][]source}{}{}{}}
\makeatother
\crefname{originalsection}{original section}{original sections}
\Crefname{originalsection}{Original section}{Original sections}
\makeatletter
\newlabel{more-algebra-section-fitting-ideals}{{more-algebra-section-fitting-ideals}{source}{}{sourcepointer-5}{}}
\newlabel{more-algebra-section-fitting-ideals@cref}{{[originalsection][0][]more-algebra-section-fitting-ideals}{[1][0][]source}{}{}{}}
\makeatother
\crefname{originallemma}{original lemma}{original lemmas}
\Crefname{originallemma}{Original lemma}{Original lemmas}
\makeatletter
\newlabel{morphisms-lemma-closed-immersion-finite-presentation}{{084Q}{source}{}{sourcepointer-6}{}}
\newlabel{morphisms-lemma-closed-immersion-finite-presentation@cref}{{[originallemma][0][]084Q}{[1][0][]source}{}{}{}}
\makeatother
\crefname{originallemma}{original lemma}{original lemmas}
\Crefname{originallemma}{Original lemma}{Original lemmas}
\makeatletter
\newlabel{properties-lemma-quasi-coherent-finite-type-ideals}{{properties-lemma-quasi-coherent-finite-type-ideals}{source}{}{sourcepointer-7}{}}
\newlabel{properties-lemma-quasi-coherent-finite-type-ideals@cref}{{[originallemma][0][]properties-lemma-quasi-coherent-finite-type-ideals}{[1][0][]source}{}{}{}}
\makeatother
\title{The Fitting ideals of a finite type module represent the loci where it is generated by r sections and where it is locally free of rank r}
\author{The Stacks Project Authors}
\date{stacks-project:divisors.tex:local-snapshot-20261012; GNU FDL 1.2 or later}
\begin{document}
\maketitle
\noindent\textit{Editorial note.} All delivered mathematical text is the author's own, copied from the pinned Stacks Project chapter file pinned-original-source.tex; no statement or proof step was rewritten, reordered or removed.\par\medskip
\noindent\textit{Editorial note.} Following the multi-claim convention for this source, source\_url is the source-level URL https://stacks.math.columbia.edu/ and the source label lemma-locally-free-rank-r-pullback is carried in claim\_key and here rather than in a per-claim tag URL.\par\medskip
\noindent\textit{Editorial note.} The preamble is the chapter's shared Stacks Project preamble, retained as sources/stacks-divisors/retained-preamble.tex: a line-for-line excerpt of the shared Stacks preamble staged with this snapshot. No macro was invented or redefined, and this card adds no preamble text of its own.\par\medskip
\noindent\textit{Editorial note.} Statement span 1383-1422 and proof span 1424-1456 come from the halfA pool card operations/pool-cards/stacks-divisors-halfA.json and were re-read in the pinned chapter file before this card was written.\par\medskip
\noindent\textit{Editorial note.} The statement contains a large displayed piecewise functor definition with a matrix-style layout; it is copied verbatim from the counted statement span.\par\medskip
\noindent\textit{Editorial note.} 7 cross-references occur inside the counted spans; they are listed in external\_references as navigation pointers only.\par\medskip
\noindent\textit{Allowed standard prerequisites (editorial).}\begin{itemize}
\item Fitting ideals of a finite type module and their behaviour over affine opens, as cited from the more-algebra chapter, including that over an affine open the n-th Fitting ideal becomes the whole ring for large n. Cited at proof lines 1425-1427. Gives part (1), the emptiness of the intersection of the closed subschemes defined by the Fitting ideals.
\item Base change of Fitting ideals and the same-chapter criteria for a module to be locally generated by r sections, respectively to be finite locally free of rank r, in terms of the r-th and the (r-1)-st Fitting ideal. Cited at proof lines 1430-1434 and 1439-1444. Converts the functorial statements (2) and (3) into statements about the vanishing of the pullback of a Fitting ideal.
\item For a finitely presented module the Fitting ideals are of finite type, so the closed immersions they cut out are of finite presentation, and a retrocompact open in the image of such an immersion is of finite presentation. Cited at proof lines 1448-1455. Gives the final finite presentation assertions about the three morphisms.
\end{itemize}\medskip
\ifdefined\hypertarget\hypertarget{sourcepointer-1}{}\fi
\noindent\textit{Source reference.} Original lemma 0C3D: focused-seeds/AGP/sources/stacks-snapshot-20261011/tags.txt, source lines 15188--15188. Same-chapter reference used in the retained author text at source line(s) 1432, 1440. It is supplied only as a navigation pointer so that the retained text compiles; the counted item delivers the author's own statement and the author's complete proof, and no part of the referenced item is copied in or substituted for the counted argument.\par\medskip
\ifdefined\hypertarget\hypertarget{sourcepointer-2}{}\fi
\noindent\textit{Source reference.} Original lemma 0C3G: focused-seeds/AGP/sources/stacks-snapshot-20261011/tags.txt, source lines 15191--15191. Same-chapter reference used in the retained author text at source line(s) 1441, 9362, 9367, 9372. It is supplied only as a navigation pointer so that the retained text compiles; the counted item delivers the author's own statement and the author's complete proof, and no part of the referenced item is copied in or substituted for the counted argument.\par\medskip
\ifdefined\hypertarget\hypertarget{sourcepointer-3}{}\fi
\noindent\textit{Source reference.} Original lemma 0C3F: focused-seeds/AGP/sources/stacks-snapshot-20261011/tags.txt, source lines 15190--15190. Same-chapter reference used in the retained author text at source line(s) 1433, 9317. It is supplied only as a navigation pointer so that the retained text compiles; the counted item delivers the author's own statement and the author's complete proof, and no part of the referenced item is copied in or substituted for the counted argument.\par\medskip
\ifdefined\hypertarget\hypertarget{sourcepointer-4}{}\fi
\noindent\textit{Source reference.} Original lemma 0C3E: focused-seeds/AGP/sources/stacks-snapshot-20261011/tags.txt, source lines 15189--15189. Same-chapter reference used in the retained author text at source line(s) 1450, 9360. It is supplied only as a navigation pointer so that the retained text compiles; the counted item delivers the author's own statement and the author's complete proof, and no part of the referenced item is copied in or substituted for the counted argument.\par\medskip
\ifdefined\hypertarget\hypertarget{sourcepointer-5}{}\fi
\noindent\textit{Source reference.} Original section more-algebra-section-fitting-ideals: sources/stacks-divisors/pinned-original-source.tex, source lines 1251--1251. Same-chapter reference used in the retained author text at source line(s) 1251, 1428. It is supplied only as a navigation pointer so that the retained text compiles; the counted item delivers the author's own statement and the author's complete proof, and no part of the referenced item is copied in or substituted for the counted argument.\par\medskip
\ifdefined\hypertarget\hypertarget{sourcepointer-6}{}\fi
\noindent\textit{Source reference.} Original lemma 084Q: focused-seeds/AGP/sources/stacks-snapshot-20261011/tags.txt, source lines 10330--10330. Same-chapter reference used in the retained author text at source line(s) 1451, 5010, 5119, 8124, 8679, 9099. It is supplied only as a navigation pointer so that the retained text compiles; the counted item delivers the author's own statement and the author's complete proof, and no part of the referenced item is copied in or substituted for the counted argument.\par\medskip
\ifdefined\hypertarget\hypertarget{sourcepointer-7}{}\fi
\noindent\textit{Source reference.} Original lemma properties-lemma-quasi-coherent-finite-type-ideals: sources/stacks-divisors/pinned-original-source.tex, source lines 1453--1453. Same-chapter reference used in the retained author text at source line(s) 1453, 8683, 9187. It is supplied only as a navigation pointer so that the retained text compiles; the counted item delivers the author's own statement and the author's complete proof, and no part of the referenced item is copied in or substituted for the counted argument.\par\medskip
\begin{lemma}
\label{lemma-locally-free-rank-r-pullback}
Let $S$ be a scheme. Let $\mathcal{F}$ be a finite type, quasi-coherent
$\mathcal{O}_S$-module. The closed subschemes
$$
S = Z_{-1} \supset Z_0 \supset Z_1 \supset Z_2 \ldots
$$
defined by the Fitting ideals of $\mathcal{F}$ have the following
properties
\begin{enumerate}
\item The intersection $\bigcap Z_r$ is empty.
\item The functor $(\Sch/S)^{opp} \to \textit{Sets}$ defined by the rule
$$
T \longmapsto
\left\{
\begin{matrix}
\{*\} & \text{if }\mathcal{F}_T\text{ is locally generated by }
\leq r\text{ sections} \\
\emptyset & \text{otherwise}
\end{matrix}
\right.
$$
is representable by the open subscheme $S \setminus Z_r$.
\item The functor $F_r : (\Sch/S)^{opp} \to \textit{Sets}$ defined by the rule
$$
T \longmapsto
\left\{
\begin{matrix}
\{*\} & \text{if }\mathcal{F}_T\text{ locally free rank }r\\
\emptyset & \text{otherwise}
\end{matrix}
\right.
$$
is representable by the locally closed subscheme $Z_{r - 1} \setminus Z_r$
of $S$.
\end{enumerate}
If $\mathcal{F}$ is of finite presentation, then
$Z_r \to S$, $S \setminus Z_r \to S$, and $Z_{r - 1} \setminus Z_r \to S$
are of finite presentation.
\end{lemma}

\begin{proof}
Part (1) is true because over every affine open $U$ there is an integer $n$
such that $\text{Fit}_n(\mathcal{F})|_U = \mathcal{O}_U$. Namely, we can
take $n$ to be the number of generators of $\mathcal{F}$ over $U$, see
More on Algebra, Section \ref{more-algebra-section-fitting-ideals}.

\medskip\noindent
For any morphism $g : T \to S$ we see from
Lemmas \ref{lemma-base-change-fitting-ideal} and
\ref{lemma-fitting-ideal-generate-locally}
that $\mathcal{F}_T$ is locally generated by $\leq r$ sections if and only if
$\text{Fit}_r(\mathcal{F}) \cdot \mathcal{O}_T = \mathcal{O}_T$.
This proves (2).

\medskip\noindent
For any morphism $g : T \to S$ we see from
Lemmas \ref{lemma-base-change-fitting-ideal} and
\ref{lemma-fitting-ideal-finite-locally-free}
that $\mathcal{F}_T$ is free of rank $r$ if and only if
$\text{Fit}_r(\mathcal{F}) \cdot \mathcal{O}_T = \mathcal{O}_T$ and
$\text{Fit}_{r - 1}(\mathcal{F}) \cdot \mathcal{O}_T = 0$.
This proves (3).

\medskip\noindent
Assume $\mathcal{F}$ is of finite presentation. Then each of the morphisms
$Z_r \to S$ is of finite presentation as $\text{Fit}_r(\mathcal{F})$
is of finite type (Lemma \ref{lemma-fitting-ideal-of-finitely-presented} and
Morphisms, Lemma \ref{morphisms-lemma-closed-immersion-finite-presentation}).
This implies that $Z_{r - 1} \setminus Z_r$ is a retrocompact open in $Z_r$
(Properties, Lemma \ref{properties-lemma-quasi-coherent-finite-type-ideals})
and hence the morphism $Z_{r - 1} \setminus Z_r \to Z_r$
is of finite presentation as well.
\end{proof}

\end{document}
